% Created 2026-02-19 Thu 21:06
% Intended LaTeX compiler: pdflatex
\documentclass[hyperref={pdfpagelabels=false},aspectratio=169]{beamer}
\usepackage[utf8]{inputenc}
\usepackage[T1]{fontenc}
\usepackage{graphicx}
\usepackage{longtable}
\usepackage{wrapfig}
\usepackage{rotating}
\usepackage[normalem]{ulem}
\usepackage{amsmath}
\usepackage{amssymb}
\usepackage{capt-of}
\usepackage{hyperref}
\usetheme{Madrid}
\usecolortheme{beaver}
\useinnertheme{rounded}
\useoutertheme{tree}
\author{Cristian A. Marocico}
\date{\today}
\title{Machine Learning in Python}
\subtitle{Neural Networks II: Deep Learning, Transformers, and GenAI}
\author[Marocico, Tatar]{Cristian A. Marocico, A. Emin Tatar}
\institute[CIT]{Center for Information Technology\\University of Groningen}
\RequirePackage{pgfcore}
\usepackage{amsfonts}
\usepackage{amsmath}
\usepackage{amssymb}
\usepackage{blkarray}
\usepackage{eurosym}
\usepackage[english]{babel}
\usepackage{xcolor,colortbl}
\usepackage[utf8]{inputenc}
\usepackage[OT1]{fontenc}
\usepackage{multirow}
\usepackage{listings}
\usepackage{fancyvrb}
\usepackage{tikz}
\usepackage{graphicx}
\usepackage[makeroom]{cancel}
\usepackage{bbm}
\newcommand{\infoTitle}[1]{\renewcommand{\givenTitle}{#1}}
\newcommand{\givenTitle}{Info}
\newenvironment{warning}[1][Info]{%
\infoTitle{#1}
\setbeamercolor{block title}{fg=white,bg=red!100!white}%
\setbeamercolor{block body}{bg=red!10!white}
\begin{block}{\givenTitle}}{
\end{block}}
\DeclareMathOperator*{\argmax}{argmax}
\setbeamercovered{transparent}
\usecolortheme{beaver}
\RequirePackage{pgfcore}
\setbeamercovered{transparent=1}
\mode<presentation>{
\usecolortheme{beaver}
\definecolor{rugcolor}{rgb}{0.8,0,0}
\definecolor{darkblue}{rgb}{0.13,0.29,0.53}
\definecolor{darkgreen}{rgb}{0.0,0.43,0.0}
\definecolor{darkyellow}{rgb}{0.0,0.43,0.43}
\definecolor{codegreen}{rgb}{0,0.6,0}
\definecolor{codegray}{rgb}{0.5,0.5,0.5}
\definecolor{codepurple}{rgb}{0.58,0,0.82}
\definecolor{backcolour}{rgb}{0.95,0.95,0.92}
\definecolor{Gray}{gray}{0.85}
\beamertemplatenavigationsymbolsempty
\setbeamercolor{item}{fg=rugcolor!80!black}
\setbeamercolor{title}{fg=rugcolor!80!black}
\setbeamercolor{frametitle}{fg=rugcolor!80!black, bg=black!10!white}
% Colors for 'definition' environment
\setbeamercolor*{block title}{fg=white, bg=darkblue}
\setbeamercolor*{block body}{bg=normal text.bg!85!darkblue}
% Color for the 'question' environment
\setbeamercolor*{block title question}{fg=white, bg=darkyellow}
\setbeamercolor*{block body question}{bg=normal text.bg!85!darkyellow}
\setbeamercolor*{palette tertiary}{bg=rugcolor,fg=white}
%HEADER WITH HIGHLIGHTED SECTION NAMES (optional)
\useheadtemplate{%
\vbox{%
%			\vskip1.2pt
%			\pgfuseimage{logo}
%			\vskip1.2pt
\tinycolouredline{rugcolor}{
\color{white}{
% comment the following line if you don't want the section names lines
%					to appear on top
\insertsectionnavigationhorizontal{\paperwidth}{}{\hskip0pt
plus1filll}
%\pgfuseimage{logored}
}
}
%    \tinycolouredline{rugcolor}
{\color{white}{
%% \insertsectionnavigationhorizontal{\paperwidth}{}{
%                \hskip0pt \hfill}
}}
}
}
%FOOTER WITH AUTHOR NAME(S), PAPER TITLE (ABBREVIATED IF SPECIFIED BY \title),
% AND PAGE COUNTER (optional)
%	\usefoottemplate{%
%		\vbox{%
%			\tinycolouredline{rugcolor}{
%				\color{white}{
%					{\insertshortauthor} \hfill \insertshortsubtitle \hfill
%					%\insertdate \hfill%
%					\textsc{\insertframenumber/\inserttotalframenumber}
%         		}
%         	}
%		}
%	}
}
\newtheorem*{props}{Properties}
\newtheorem*{prop}{Property}
\newtheorem*{notation}{Notation}
\newtheorem*{terminology}{Terminology}
\newcolumntype{a}{>{\columncolor{Gray}}c}
\DeclareMathOperator*{\argmin}{argmin}
\newcommand{\Var}{\mathbb{V}\mathrm{ar}}
\newcommand{\Corr}{\mathbb{C}\mathrm{orr}}
\newcommand{\Cov}{\mathbb{C}\mathrm{ov}}
\newcommand{\Expt}{\mathbb{E}}
\newcommand{\NorDist}{\mathcal{N}}
\newcommand{\ExpDist}{\mathcal{E}\mathrm{xp}}
\newcommand{\GammaDist}{\mathcal{G}\mathrm{amma}}
\newcommand{\BetaDist}{\mathcal{B}\mathrm{eta}}
\newcounter{listCounter}
\newenvironment{question}{%
\setbeamercolor{block title}{bg=orange!70!white,fg=white}
\setbeamercolor{block body}{bg=yellow!10!white}
\begin{block}{Question}
}{%
\end{block}
}
\lstdefinestyle{mystyle}{
language=R,
backgroundcolor=\color{backcolour},
commentstyle=\color{codegreen},
keywordstyle=\color{darkblue}\bfseries,
numberstyle=\tiny\color{codegray},
stringstyle=\color{codepurple},
basicstyle=\ttfamily,
breakatwhitespace=true,
breaklines=true,
%	captionpos=none,
columns=fixed,
keepspaces=true,
numbers=none,
numbersep=5pt,
showspaces=false,
showstringspaces=false,
showtabs=false,
tabsize=2,
basewidth=1.5em,
escapeinside={<@}{@>}
}
\setbeamertemplate{caption}{\raggedright\insertcaption\par}
\setbeamertemplate{blocks}[rounded][shadow=true]
\renewenvironment{definition}[1][Definition]{%
\setbeamercolor{block title}{fg=white,bg=darkblue}
\setbeamercolor{block body}{fg=black,bg=normal text.bg!85!darkblue}
\begin{block}{#1\hfill \footnotesize{Definition}}
\vspace*{-5pt}
}{%
\end{block}
}
\renewenvironment{example}[1][Example]{%
% Color for the 'example' environment
\setbeamercolor*{block title}{fg=white, bg=darkgreen}
\setbeamercolor*{block body}{bg=normal text.bg!85!darkgreen}
\begin{block}{#1\hfill \footnotesize{Example}}
\vspace*{-5pt}
}{%
\end{block}
}
\renewenvironment{theorem}[1][Theorem]{%
\setbeamercolor*{block title}{fg=white, bg=darkyellow}
\setbeamercolor*{block body}{bg=normal text.bg!85!darkyellow}
\begin{block}{#1\hfill \footnotesize{Theorem}}
\vspace*{-5pt}
}{%
\end{block}
}
\date[Feb 20\textsuperscript{th} 2026]{Friday, February 20\textsuperscript{th} 2026}
\usepackage{mathtools}
\setbeamercovered{transparent=0}
\usepackage[timeinterval=60]{tdclock}
\hypersetup{
 pdfauthor={Cristian A. Marocico},
 pdftitle={Machine Learning in Python},
 pdfkeywords={},
 pdfsubject={},
 pdfcreator={Emacs 30.2 (Org mode 9.7.11)}, 
 pdflang={English}}
\begin{document}

\maketitle
\begin{frame}{Outline}
\setcounter{tocdepth}{1}
\tableofcontents
\end{frame}

\section{Convolutional Neural Networks (CNNs)}
\label{sec:org70519a5}
\begin{frame}[label={sec:orge37f4f6}]{Why CNNs?}
\begin{definition}[The Problem with MLPs for Images]\label{sec:org9ac777c}
\pause
Standard MLPs flatten images into vectors, causing two major problems:
\begin{itemize}[<+->]
\item \alert{Spatial Structure Lost}: A nose is defined by being \emph{above} a mouth, not just by pixel values.
\item \alert{Parameter Explosion}: A \(1000 \times 1000\) image connected to 100 neurons requires \alert{100 million weights}.
\end{itemize}
\end{definition}
\end{frame}
\begin{frame}[label={sec:org8c54d23}]{The Convolution Operation}
\begin{example}[Filters and Feature Maps]\label{sec:org0b6d308}
\pause
Instead of learning weights for every pixel, CNNs learn small \alert{Filters} (Kernels):
\begin{itemize}[<+->]
\item A filter might be \(3 \times 3\) pixels.
\item It slides (convolves) over the image, computing dot products.
\item This creates a \alert{Feature Map} showing where patterns appear.
\item Early filters detect edges; deeper filters detect complex shapes.
\end{itemize}
\end{example}
\end{frame}
\begin{frame}[label={sec:org0769352}]{Controlling Dimensions}
\begin{block}{Padding and Stride}
\pause
As filters slide, image dimensions change:
\begin{itemize}[<+->]
\item \alert{Valid Convolution}: No padding. Output shrinks (\(5 \times 5 \to 3 \times 3\)).
\item \alert{Padding}: Add zeros around border to maintain size.
\item \alert{Stride}: How many pixels to jump. Stride=2 usually reduces dimensions by about half.
\item \alert{Formula}: \(O = \frac{I - K + 2P}{S} + 1\)
\end{itemize}
\end{block}
\end{frame}
\begin{frame}[label={sec:orge0f3f6b}]{Pooling Layers}
\begin{example}[MaxPool2d]\label{sec:org388c6ac}
\pause
Pooling reduces spatial size to control overfitting and computation:
\begin{itemize}[<+->]
\item \alert{Max Pooling}: Takes the largest value in a window.
\item Asks: "Is there a feature in this patch? Yes/No."
\item Provides translation invariance (feature can shift slightly).
\item Typical: \(2 \times 2\) pool with stride 2 halves dimensions.
\end{itemize}
\end{example}
\end{frame}
\begin{frame}[label={sec:org43db836}]{Famous Architectures}
\begin{block}{VGG and ResNet}
\pause
Modern Deep Learning is built on famous architectures:
\begin{itemize}[<+->]
\item \alert{VGG (2014)}: Simplicity. Stacks many small \(3 \times 3\) filters. Very deep.
\item \alert{ResNet (2015)}: Introduced \alert{Skip Connections} (Residuals).
\item Skip connections let gradients "skip" layers to avoid vanishing.
\item Enables training networks with 100+ layers.
\end{itemize}
\end{block}
\end{frame}
\section{From RNNs to Transformers}
\label{sec:orgd6c93a0}
\begin{frame}[label={sec:org14d23b8}]{Why This Matters}
\begin{definition}[Sequence Data Is Everywhere]\label{sec:orgeb04f7e}
\pause
Many research problems are sequential:
\begin{itemize}[<+->]
\item Clinical notes and patient timelines
\item Instrument logs and time-ordered measurements
\item Reports, interviews, and long documents
\item Key question: \emph{which earlier tokens matter right now?}
\end{itemize}
\end{definition}
\end{frame}
\begin{frame}[label={sec:org7576a41}]{Running Example}
\begin{example}[Clinical Sentence]\label{sec:org0e3f8c8}
\pause
\emph{"Patient admitted three days ago with acute chest pain, treated with aspirin, now reports reduced discomfort."}
\pause
When interpreting \alert{"discomfort"}, the model should retrieve context from:
\begin{itemize}[<+->]
\item "chest pain"
\item "aspirin"
\item other relevant earlier tokens
\end{itemize}
\end{example}
\end{frame}
\begin{frame}[label={sec:org2772e15}]{RNNs: First Main Solution}
\begin{block}{Sequential Hidden-State Memory}
\pause
RNN update:
$$h_t = \tanh(Wx_t + Uh_{t-1})$$
\begin{itemize}[<+->]
\item \(x_t\): current token representation
\item \(h_{t-1}\): previous memory
\item \(h_t\): updated memory
\item Strength: natural for ordered data
\item Limits: sequential bottleneck + fading long-range memory
\item LSTMs improved memory with gates, but remained sequential
\end{itemize}
\end{block}
\end{frame}
\begin{frame}[label={sec:org6343a19}]{Attention as Weighted Lookup}
\begin{example}[Query-Specific Relevance Weights]\label{sec:org07ec361}
\pause
Attention replaces "single compressed memory" with dynamic lookup.
\begin{itemize}[<+->]
\item Weights are non-negative and sum to 1
\item Different query tokens produce different weight patterns
\item For query token "today", weights are distributed across ["Patient", "reports", "chest", "pain", "today"]
\item Self-weight can be < 1 because context tokens also contribute
\end{itemize}
\end{example}
\end{frame}
\begin{frame}[label={sec:org5cc7223}]{Transformers: Attention-Centered Architecture}
\begin{definition}["Attention Is All You Need"]\label{sec:orgcb3de4a}
\pause
Transformers remove recurrence and rely on attention:
\begin{center}
\begin{tabular}{lll}
Feature & RNN / LSTM & Transformer\\
\hline
Processing & Sequential & Parallel in training\\
Context & Compressed in one hidden state & Direct weighted access across positions\\
Long distance & Degrades with sequence length & One attention step away (subject to masking rule)\\
\end{tabular}
\end{center}
\end{definition}
\end{frame}
\begin{frame}[label={sec:orgc4d3363}]{Core Formula and Transformer Block}
\begin{block}{Q, K, V and Scaled Dot Product}
\pause
$$\text{Attention}(Q, K, V) = \text{softmax}\left(\frac{QK^T}{\sqrt{d_k}}\right)V$$
\begin{itemize}[<+->]
\item Query: what this position is looking for
\item Key: what each position offers
\item Value: information to aggregate
\item \(\sqrt{d_k}\) scaling stabilizes softmax/gradients
\item Block components: self-attention, feed-forward, residual, layer norm, positional encoding
\end{itemize}
\end{block}
\end{frame}
\begin{frame}[label={sec:orgf47baab}]{Multi-Head and Model Families}
\begin{example}[Parallel Attention Heads]\label{sec:org1821832}
\pause
Multi-head attention runs several Q/K/V projections in parallel.
\begin{itemize}[<+->]
\item Different heads can focus on different relations (syntax, semantics, locality)
\item Outputs are concatenated and projected back to model dimension
\item Encoder-only: BERT (understanding/classification)
\item Decoder-only: GPT (generation with causal masking)
\item Encoder-decoder: T5 (translation/summarization)
\end{itemize}
\end{example}
\end{frame}
\begin{frame}[label={sec:orga77f255}]{Why GPT Is Left-to-Right}
\begin{block}{Causal Masking}
\pause
For token \(t\), attention is allowed only to tokens \(\le t\).
\begin{itemize}[<+->]
\item Future positions are masked out
\item Training remains parallel across positions (with masking)
\item Generation remains token-by-token
\item Early tokens have less context than later tokens
\end{itemize}
\end{block}
\end{frame}
\begin{frame}[label={sec:orgb107b84}]{Practical Research Workflow}
\begin{example}[Pretrained-First Strategy]\label{sec:orga32f035}
\pause
Most researchers do not train transformers from scratch:
\begin{itemize}[<+->]
\item Choose a pretrained model for the task/domain
\item Start with prompting and retrieval
\item Evaluate on domain-specific examples
\item Fine-tune only if needed
\item Keep expert review in high-stakes settings
\end{itemize}
\end{example}
\end{frame}
\section{Large Language Models (LLMs)}
\label{sec:org0717edf}
\begin{frame}[label={sec:orgd621ae1}]{Evolution of LLMs}
\begin{definition}[Scale as the Main Driver]\label{sec:org6c1e75b}
\pause
LLM progress accelerated with scale (parameters, data, compute):
\begin{center}
\begin{tabular}{lrlll}
Model & Year & Parameters & Training Data & Notable Point\\
\hline
\alert{BERT} & 2018 & 340M & 16GB & Strong "understanding" performance\\
\alert{GPT-2} & 2019 & 1.5B & 40GB & High-quality open-ended generation\\
\alert{GPT-3} & 2020 & 175B & 570GB & Emergent few-shot behavior\\
\alert{Llama-2} & 2023 & 70B & 2TB & Open-weight ecosystem growth\\
\alert{GPT-4} & 2023 & Undisclosed & Unknown & Strong multimodal capabilities\\
\end{tabular}
\end{center}
\end{definition}
\end{frame}
\begin{frame}[label={sec:org2967028}]{Architecture Families}
\begin{block}{Encoder, Decoder, Encoder-Decoder}
\pause
\begin{itemize}[<+->]
\item \alert{BERT (Encoder-only)}: understanding tasks (classification, QA)
\item \alert{GPT (Decoder-only)}: generation tasks (next-token prediction)
\item \alert{T5 (Encoder-Decoder)}: text-to-text tasks (translation, summarization)
\item \alert{Scaling hypothesis}: performance improves predictably with more parameters, data, and compute
\end{itemize}
\end{block}
\end{frame}
\begin{frame}[label={sec:orgdb25937}]{The LLM Lifecycle}
\begin{example}[Pre-training, Fine-tuning, Inference]\label{sec:org4356ec1}
\pause
\begin{itemize}[<+->]
\item \alert{Pre-training}: The expensive part. Read the internet (TB of text), learn to predict next token. Learns grammar, facts, reasoning.
\item \alert{Fine-tuning}: Specialization.
\begin{itemize}
\item \alert{SFT}: Train on Q\&A pairs to make it a chatbot.
\item \alert{RLHF}: Align to be helpful and harmless.
\end{itemize}
\item \alert{Inference}: Using the model. Feed prompt, generate response.
\end{itemize}
\end{example}
\end{frame}
\begin{frame}[label={sec:orgce35b27}]{Tokenization and Context Windows}
\begin{example}[BPE and Subword Units]\label{sec:orgb995ac7}
\pause
LLMs don't see words; they see \alert{Tokens}:
\begin{itemize}[<+->]
\item A token can be a word ("apple"), part of a word ("ing"), or a character.
\item \alert{Rough rule}: 1000 tokens \(\approx\) 750 words.
\item BPE/subword tokenization handles rare or unseen words by splitting them.
\item \alert{Context Window}: Limit on how much text the model can use at once (model/version dependent).
\end{itemize}
\end{example}
\end{frame}
\begin{frame}[label={sec:orgb8769d1}]{Next-Token Prediction Mechanics}
\begin{block}{From Logits to Text}
\pause
At each decoding step, the model outputs logits \(z \in \mathbb{R}^{|V|}\) over the vocabulary.
\pause

Probabilities are computed as:
\[
p_i = \frac{e^{z_i}}{\sum_j e^{z_j}}
\]
\pause
\begin{itemize}[<+->]
\item This is a full distribution over all candidate next tokens.
\item Decoding selects/samples one token from this distribution.
\item The selected token is appended, and the process repeats autoregressively.
\end{itemize}
\end{block}
\end{frame}
\begin{frame}[label={sec:org42e5ea8}]{Sampling Controls}
\begin{example}[Temperature, Top-k, Top-p]\label{sec:org55a51a2}
\pause
Generation quality/style depends on sampling strategy:
\begin{center}
\begin{tabular}{llr}
Parameter & Effect & Typical\\
\hline
\alert{Temperature} & Flattens/sharpens token probabilities & 0.7\\
\alert{Top-k} & Restricts choices to k most likely tokens & 50\\
\alert{Top-p} & Restricts to smallest set with cumulative prob \(\ge p\) & 0.9\\
\end{tabular}
\end{center}
\pause
\begin{itemize}[<+->]
\item Low temperature: focused/deterministic
\item High temperature: diverse/chaotic
\item Top-k/top-p reduce low-probability tail sampling
\end{itemize}
\end{example}
\end{frame}
\begin{frame}[label={sec:org8171ca6}]{Prompt Engineering}
\begin{example}[Guiding the Predictor]\label{sec:org150d137}
\pause
LLMs are "next token predictors." They are sensitive to \emph{how} you ask:
\pause
\begin{itemize}[<+->]
\item \alert{Zero-Shot}: Just asking. "Translate to Spanish: Hello"
\item \alert{Few-Shot}: Give examples to induce pattern-following behavior.
\item \alert{Chain of Thought}: "Let's think step by step." Can help some multi-step tasks (model/task dependent).
\end{itemize}
\end{example}
\end{frame}
\section{Generative AI and Modern Applications}
\label{sec:org5157b4f}
\begin{frame}[label={sec:org9ce0a6c}]{Text Generation}
\begin{definition}[It's Just Statistics]\label{sec:org586c2af}
\pause
At its core, a Generative LLM is a \alert{Next-Token Prediction Engine}:
\pause
Given "The sky is", it calculates probabilities:
\begin{itemize}[<+->]
\item "blue": 80\%
\item "gray": 15\%
\item "falling": 1\%
\item It samples from this distribution to generate text.
\end{itemize}
\end{definition}
\end{frame}
\begin{frame}[label={sec:orgcebf702}]{Retrieval-Augmented Generation (RAG)}
\begin{example}[Solving LLM Limitations]\label{sec:org97c7202}
\pause
LLMs have two major flaws:
\begin{itemize}[<+->]
\item \alert{Hallucination}: They confidently make things up.
\item \alert{Cut-off Date}: They don't know current events or private data.
\item One practical mitigation is \alert{RAG}, which connects LLMs to external knowledge:
\begin{enumerate}
\item \alert{Retrieve}: Search database for relevant documents.
\item \alert{Augment}: Paste documents into prompt context.
\item \alert{Generate}: LLM answers using provided context.
\end{enumerate}
\end{itemize}
\end{example}
\end{frame}
\begin{frame}[label={sec:org9907df3}]{Diffusion Models}
\begin{block}{Image Generation (Stable Diffusion, DALL-E)}
\pause
Image generation works differently from LLMs, using \alert{Diffusion}:
\begin{itemize}[<+->]
\item \alert{Forward Process (Training)}: Take clear image, slowly add noise until pure static. Model learns destruction pattern.
\item \alert{Reverse Process (Generation)}: Start with noise. Ask: "If this was a cat, how would you remove noise?" Repeat 50 times.
\item Result: Clear image emerges from noise.
\end{itemize}
\end{block}
\end{frame}
\begin{frame}[label={sec:org7f503aa}]{Ethical Considerations}
\begin{example}[Bias, Hallucination, and Copyright]\label{sec:org87b58a3}
\pause
With great power comes great responsibility:
\begin{itemize}[<+->]
\item \alert{Bias}: Models trained on internet data reflect stereotypes. "Doctor" \(\to\) men, "Nurse" \(\to\) women.
\item \alert{Hallucination}: LLMs optimize for plausibility, not truth. Dangerous for medical/legal advice.
\item \alert{Copyright}: Image models trained on artists' work. Code models on GitHub. Legal framework still evolving.
\item \alert{Mitigation}: Data curation, RLHF, human oversight.
\end{itemize}
\end{example}
\end{frame}
\section{Time-Series Forecasting with Feature-Based ML}
\label{sec:org6311b41}
\begin{frame}[label={sec:org3470a35}]{Forecasting Problem Setup}
\begin{definition}[Formal Objective]\label{sec:org6327625}
\pause
For a univariate series \(\{y_t\}_{t=1}^T\), one-step forecasting is:
\[
\hat{y}_{t+1|t} = f_{\theta}(\phi_t)
\]
where \(\phi_t\) uses only information available at time \(t\).
\pause
Training is empirical risk minimization:
\[
\hat{\theta} = \arg\min_{\theta}\frac{1}{N}\sum_{t \in \mathcal{T}_{train}} \mathcal{L}\left(y_{t+1}, f_{\theta}(\phi_t)\right)
\]
\end{definition}
\end{frame}
\begin{frame}[label={sec:org6898795}]{Why Time Series Is Not i.i.d.}
\begin{block}{Temporal Dependence and Structure}
\pause
Time series typically combines:
\[
y_t = T_t + S_t + R_t
\]
\pause
\begin{itemize}[<+->]
\item \(T_t\): trend (long-term movement)
\item \(S_t\): seasonality (repeating calendar pattern)
\item \(R_t\): residual/noise
\end{itemize}
\end{block}
\end{frame}
\begin{frame}[label={sec:orgdf0cb2f}]{Why Random Splits Fail}
\begin{example}[Leakage Through Time Mixing]\label{sec:org10432bb}
\pause
Information sets are ordered:
\[
\mathcal{F}_{t_1} \subset \mathcal{F}_{t_2}, \quad t_1 < t_2
\]
\pause
Random shuffling violates this causal flow and can leak future-conditioned patterns into training.
\pause
\begin{itemize}[<+->]
\item Correct: train on past, validate on future block.
\item Incorrect: random train/test split across timestamps.
\end{itemize}
\end{example}
\end{frame}
\begin{frame}[label={sec:orga612268}]{Strong Baselines First}
\begin{definition}[Naive and Seasonal Naive]\label{sec:org62e79d9}
\pause
\[
\hat{y}_{t + 1}^{naive}=y_t,\qquad
\hat{y}_{t + 1}^{seasonal}=y_{t + 1 - s}
\]
\pause
for monthly data, \(s=12\).
\pause
\begin{itemize}[<+->]
\item Random walk \(y_t=y_{t - 1} + \varepsilon_t \rightarrow\) naive is optimal (MSE).
\item Seasonal random walk \(y_t=y_{t-s}+\varepsilon_t \rightarrow\) seasonal naive is optimal.
\item If ML cannot beat these, complexity is not justified.
\end{itemize}
\end{definition}
\end{frame}
\begin{frame}[label={sec:org9d7df7d}]{Feature Engineering with Temporal Constraints}
\begin{block}{Building \(\phi_t\) Safely}
\pause
Example feature vector:
\[
\phi_t = [y_{t-1}, y_{t-2}, y_{t-12}, \mu_{t-1}^{(3)}, \mu_{t-1}^{(12)},
\sin(2\pi m_t/12), \cos(2\pi m_t/12), t]
\]
with rolling mean
\[
\mu_{t-1}^{(w)} = \frac{1}{w}\sum_{j=1}^{w} y_{t-j}
\]
\begin{itemize}[<+->]
\item Shift first, then roll (prevents target leakage).
\item Cyclical encoding respects December-January adjacency.
\end{itemize}
\end{block}
\end{frame}
\begin{frame}[label={sec:org56c118b}]{Model Families and Objectives}
\begin{example}[Ridge vs Boosting]\label{sec:orgfce97c4}
\pause
\begin{itemize}[<+->]
\item Ridge: stable and interpretable with correlated lags. Ridge objective:
\end{itemize}
\[
\hat{\beta}=\arg\min_{\beta}\left[\frac{1}{n}|y - X\beta|^2+\alpha|\beta|^2\right]
\]
\begin{itemize}
\item Boosting: captures nonlinear interactions. Boosting (additive model):
\end{itemize}
\[
f_M(x)=\sum_{m=1}^M \eta h_m(x),\qquad
r_i^{(m)} \approx y_i - f_{m-1}(x_i)
\]
\end{example}
\end{frame}
\begin{frame}[label={sec:org7b42126}]{Walk-Forward Backtesting}
\begin{example}[Estimating Future Performance]\label{sec:org1967bd1}
\pause
Choose ordered cutoffs \(\tau_1 < \dots < \tau_K\). For fold \(k\):
\begin{itemize}[<+->]
\item Train on \(\{1,\dots,\tau_k\}\)
\item Test on \(\{\tau_k+1,\dots,\tau_k+h\}\)
\item Compute metric \(M_k\)
\end{itemize}
\pause
Aggregate:
\[
\bar{M}=\frac{1}{K}\sum_{k=1}^K M_k
\]
This is usually more reliable than a single holdout split.
\end{example}
\end{frame}
\begin{frame}[label={sec:orgf7650b2}]{Diagnostics After Forecasting}
\begin{block}{Residual Checks}
\pause
Even with good MAE/RMSE, verify residual structure.
\begin{itemize}[<+->]
\item Residual mean near zero (bias check).
\item No strong temporal pattern in residual plots.
\item Low residual autocorrelation at small lags.
\item If autocorrelation remains, predictive signal is still unmodeled.
\end{itemize}
\end{block}
\end{frame}
\section{Summary}
\label{sec:orgd620280}
\begin{frame}[label={sec:org04a9aa1}]{Key Takeaways}
\begin{block}{Neural Networks II}
\pause
\begin{enumerate}[<+->]
\item \alert{CNNs}: Use convolutions and pooling to efficiently process images. ResNet enables very deep networks.
\item \alert{Transformers}: Replaced RNNs with parallel self-attention. The foundation of modern NLP.
\item \alert{LLMs}: Giant Transformers (BERT, GPT, Llama) that predict next tokens. Power chatbots and assistants.
\item \alert{Time-Series ML}: Formalize \(f_\theta(\phi_t)\), beat strong baselines, and validate with walk-forward backtesting.
\item \alert{Generative AI}: Text (next-token), Images (diffusion). RAG grounds LLMs in facts.
\item \alert{Ethics}: AI reflects training data. Requires human oversight and responsibility.
\end{enumerate}
\end{block}
\end{frame}
\end{document}
